% SEGMENT OLP-0689-S01
% Part: proof-theory
% Chapter: propositions-as-types
% Section: proofs-to-terms

\documentclass[../../../include/open-logic-section]{subfiles}

\begin{document}

\olfileid{int}{pty}{pt}

\olsection{\usetoken{P}{derivation} ஐ நிறுவல் சொற்களாக மாற்றுதல்}

இயல்பான வருவித்தலின் !!{proof}s ஐ நிறுவல் சொற்களாக
மாற்றுவது, அதாவது \Log{N2i} இன் !!{proof}s ஐ
\Log{tN2i} இன் !!{proof}s ஆக மொழிமாற்றுவது,
மிகவும் நேரடியான செயல். அந்த !!{derivation} இல்
ஒவ்வொரு தொடருருவின் வலப்புற !!{formula} வுக்கும்
இடப்புறம் ஒரு நிறுவல் சொல்லை எழுதினால் போதும்;
அப்போது $\Gamma \Sequent M : !A$ என்ற வடிவத்
தொடருருக்கள் கிடைக்கும். பின்னர், சூழல்~$\Gamma$ இல்
$!A$ க்கு $M$ \emph{சான்றாக அமைகிறது} என்போம்.
எந்த நிறுவல் சொல்லை எழுதுவது என்பதை
\Log{tN2i} இன் விதிகளே முழுவதும் தீர்மானிக்கின்றன.

\begin{prop}
\Log{N2i} இல் $\Gamma \Sequent !A$ என்பதற்கான
ஒவ்வொரு !!{proof}~$\delta$ வுக்கும், \Log{tN2i} இல்
$\Gamma \Sequent N : !A$ என்பதற்கான !!a{proof}~$\delta'$
பொருந்தும். நிறுவல் சொல்~$N$ ஐ $\delta$ தனித்துத்
தீர்மானிக்கிறது.
\end{prop}

% SEGMENT OLP-0689-S02
\begin{proof}
$\delta$ வின் உயரத்தின்மீது தொகுத்தறிதல் செய்வோம்.
$\delta$ என்பது $x:!A \Sequent !A$ என்ற அடிக்கோள்
என்றால், $\delta'$ என்பது $x:!A \Sequent x:!A$
என்ற அடிக்கோள். $\delta$ ஓர் அனுமான விதியில் முடிந்தால்,
தொகுத்தறிதல் கருதுகோள் அதன் ஒவ்வொரு முன்கூற்றுக்கும்
ஒரு நிறுவல் சொல்லையும், அந்த முன்கூற்றுக்குப் பொருந்திய
நிறுவல் சொல் அடங்கிய தொடருருவிற்கான
\Log{tN2i}-!!{proof}s ஐயும் தருகிறது. அதற்குப்
பொருந்திய \Log{tN2i} விதியைப் பயன்படுத்துகிறோம்;
முடிவின் நிறுவல் சொல்லை அந்த \Log{tN2i} அனுமான
விதியே தீர்மானிக்கிறது.
\end{proof}

% SEGMENT OLP-0689-S03
\begin{ex}
  $(!A \land !B) \lif (!B \land !A)$ என்பதற்கான
  \Log{N1i} இன் மிக எளிய !!{proof} ஒன்றைக் காண்போம்:
  \[
  \AxiomC{$\Discharge{!A\land !B}{x}$}
  \RightLabel{\Elim\land[2]}
  \UnaryInfC{$!B$}
  \AxiomC{$\Discharge{!A\land !B}{x}$}
  \RightLabel{\Elim\land[1]}
  \UnaryInfC{$!A$}
  \RightLabel{\Intro\land}
  \BinaryInfC{$!B \land !A$}
  \DischargeRule{\Intro\lif}{x}
  \UnaryInfC{$(!A \land !B) \lif (!B \land !A)$}
  \DisplayProof
  \]

% SEGMENT OLP-0689-S04
  \Log{N2i} இல் அது பின்வருமாறு அமையும்:
  \[
  \Axiom$x : !A\land !B \fCenter !A\land !B$
  \RightLabel{\Elim\land[2]}
  \UnaryInf$x : !A\land !B \fCenter !B$
  \Axiom$x : !A\land !B \fCenter !A\land !B$
  \RightLabel{\Elim\land[1]}
  \UnaryInf$x : !A\land !B \fCenter !A$
  \RightLabel{\Intro\land}
  \BinaryInf$x : !A\land !B \fCenter !B \land !A$
  \RightLabel{\Intro\lif}
  \UnaryInf$\fCenter (!A \land !B) \lif (!B \land !A)$
  \DisplayProof
  \]

% SEGMENT OLP-0689-S05
  நிறுவல் சொற்களை மேலிருந்து கீழாக ஒதுக்குகிறோம்.
  அடிக்கோள்களுக்குரிய நிறுவல் சொல், சூழலிலுள்ள
  அதே மாறியான~$x$. $\Elim\land[i]$ விதிகளுக்குரிய
  நிறுவல் சொற்கள் முறையே $\proj{2}{x}$ என்றும்
  $\proj{1}{x}$ என்றும் தீர்மானிக்கப்படுகின்றன.
  \Intro{\land} விதியைப் பயன்படுத்தும்போது,
  இவ்விரு நிறுவல் சொற்களையும்
  $\pair{\proj{2}{x}}{\proj{1}{x}}$ என இணைக்கிறோம்.
  இறுதியாக, \Intro{\lif} விதி ஒரு $\lambd$-சுருக்கத்தைப்
  பயன்படுத்துகிறது: அது~$x$ ஐக் கட்டுப்படுத்தி,
  $!A \land !B$ என்ற அனுமானத்தின் அடையாளம்~$x$
  என்பதையும் பதிவுசெய்கிறது. ஆகவே நிறுவல் சொல்
  $\lambd[\typeof{x}{!A \land
  !B}][\pair{\proj{2}{x}}{\proj{1}{x}}]$.
  \Log{tN2i} இல் உள்ள !!{proof} பின்வருமாறு:
  \[
  \Axiom$x : !A\land !B \fCenter x: !A\land !B$
  \RightLabel{\Elim\land[2]}
  \UnaryInf$x : !A\land !B \fCenter \proj{2}{x} : !B$
  \Axiom$x : !A\land !B \fCenter x: !A\land !B$
  \RightLabel{\Elim\land[1]}
  \UnaryInf$x : !A\land !B \fCenter \proj{1}{x} : !A$
  \RightLabel{\Intro\land}
  \BinaryInf$x : !A\land !B \fCenter \pair{\proj{2}{x}}{\proj{1}{x}} :
    !B \land !A$
  \RightLabel{\Intro\lif}
  \UnaryInf$\fCenter \lambd[\typeof{x}{!A \land
  !B}][\pair{\proj{2}{x}}{\proj{1}{x}}] : (!A \land !B) \lif (!B \land !A)$
  \DisplayProof
  \]
\end{ex}

\end{document}
